\section{Introduction}
\label{sec:introduction}

The Weisfeiler--Leman algorithm refines the colors of tuples of vertices
to test graph isomorphism. A graph's Weisfeiler--Leman dimension measures how much tuple
information is needed to identify a particular graph among all graphs.
This differs from testing whether a specified pair of graphs survives
refinement: identification quantifies over every possible comparison
graph. We determine the complexity of identification when the dimension
is part of the input.

Increasing the dimension exposes more information, but also increases
the number of tuples on which refinement operates. This makes the least
sufficient dimension a natural measure of the strength needed on an
individual graph. The decision problem studied here asks for that
threshold, rather than for the number of refinement rounds or the cost
of one particular implementation of refinement.

All graphs in this paper are finite, simple, undirected, and uncolored.
For $k\ge2$ we use the joint-update convention: a $k$-tuple starts with
its equality and adjacency type, and each refinement adjoins the
multiset, over vertices $z$, of the vectors of colors obtained by
replacing each coordinate by that same $z$. At $k=1$ we use ordinary
neighbor color refinement. Colors are shared between the graphs being
compared. Write $G\equiv_k H$ when their tuple-color histograms agree
at every round. We say that $k$-WL \emph{identifies} $G$ if
$G\equiv_k H$ implies $G\cong H$ for every graph $H$.
The \emph{Weisfeiler--Leman dimension} $\WLdim(G)$ is the least positive
integer $k$ that identifies $G$. Section~\ref{sec:wl} gives the formal
refinement rule and proves the monotonicity needed for this definition.

Our decision problem has as input a nonempty graph $G$, given by its
binary adjacency matrix, and a positive integer $k$, given in binary.
It asks whether $\WLdim(G)\le k$. Malformed inputs are rejected.

\begin{theorem}
\label{thm:main}
The following problem is EXPTIME-complete under deterministic
polynomial-time many-one reductions: given a nonempty finite simple
undirected uncolored graph $G$ by its adjacency matrix and a positive
integer $k$ in binary, decide whether $\WLdim(G)\le k$.
Here dimension one uses ordinary neighbor color refinement, and
dimensions at least two use the joint-update rule described above.
\end{theorem}

The problem remains EXPTIME-complete when $k$ is encoded in unary
(Corollary~\ref{cor:unary}). Thus the hardness does not require a
dimension whose numerical value is exponentially larger than its encoding.

\subsection{Context and contribution}

The refinement method originates in Weisfeiler and Leman's work on
graph canonization and its associated algebra~\cite{WL1968}.
The relation between Weisfeiler--Leman refinement, counting logic, and
pebble games is a central tool in finite model theory and graph
isomorphism; see Cai, F\"urer, and Immerman~\cite{CFI1992} and the
game formulation in Kiefer and Neuen~\cite{KN2019}.
Arvind, K\"obler, Rattan, and Verbitsky proved P-hardness for
dimension-one identification~\cite{AKRV2017}. Lichter, Ra\ss mann,
and Schweitzer extended fixed-dimension P-hardness to every $k\ge2$
and proved NP-hardness when the dimension is part of the input,
including simple uncolored graphs
\cite[Theorems~3--4]{LRS2025}.
Those results distinguish identification from equivalence of a given
pair. Variable-dimension pair equivalence was shown coNP-hard by
Seppelt~\cite[Theorem~5]{Seppelt2024} and independently by Lichter,
Ra\ss mann, and Schweitzer~\cite[Theorem~19]{LRS2025}.

The identification arguments of Lichter, Ra\ss mann, and Schweitzer
already analyze every equivalent mate of their colored CFI graphs, using
local recognition of gadget incidences
\cite[Lemma~8 and Claim~9]{LRS2024full}. The issue here is therefore
not the all-mates quantifier by itself. Our incidence encoding admits
an entire family of shifted constraints, and we must prove that every
shift arising from an equivalent mate lifts to one global isomorphism
on an accepting instance. This requires more than recognizing the
allowed constraint pattern one incidence piece at a time.

At the level of graph encodings, local parity constraints and a global
twist obstruction are central features of the Cai--F\"urer--Immerman
construction~\cite{CFI1992}. Compressed CFI graphs provide a related
method for retaining game obstructions in smaller representations
\cite{GLNS2025}. That work proves lower bounds on refinement rounds;
such bounds do not by themselves determine the complexity of deciding
identification. Our representation and its game properties are proved
directly below.

The companion paper \emph{Variable-dimension Weisfeiler--Leman
equivalence on general and subcubic graphs}~\cite{VariableWL} develops
a compressed computation construction for pair equivalence. We use its
short tables on binary coordinates, comparisons of images of scalar
linear forms, and potential-based extension of local assignments.
The identification-specific arguments are proved below;
in particular, pair-equivalence hardness is not taken to imply
identification hardness. The graph representation is changed so
that the structure of an arbitrary equivalent graph can be recovered
from bounded-size incidence tests.

Theorem~\ref{thm:main} completes the complexity classification of the
input-dimension identification problem. The additional mechanism is a
local-to-global reconstruction for binary linear constraints. Small
incidence pieces determine the possible shifts in any equivalent graph.
We distinguish local choices that suffice for a winning game strategy
from compatible choices that define one graph isomorphism. Both
conditions are established for an abstract family of comparisons before
the comparisons are chosen to encode a computation.

\subsection{Proof outline}

For a fixed exponential-time machine and input word, the reduction
constructs a graph $G(0)$ and a polynomially bounded dimension $r$.
The graph records binary linear constraints. Each variable bit has a
pair of vertices, and each small constraint has vertices for its legal
valuations, joined to the bit values they specify. Distinct colors for
these classes are subsequently encoded by private leaves. The notation
$G(\beta)$ denotes the same construction with constraint right-hand
sides changed from zero to a list of bits $\beta$.

The representation has three properties. First, every graph equivalent
to $G(0)$ at dimension $r$ is isomorphic to some $G(\beta)$, and such
a mate is isomorphic to $G(0)$ exactly when its shifts admit a global
solution. Second, game equivalence yields nonempty sets of vector
values indexed by binary addresses $a\in\{0,1\}^r$. Each comparison
equates just one scalar image of these sets, up to the prescribed
shift. For a particular class of shifts, these image conditions also
suffice for game equivalence. Third, a global solution exists exactly
when the comparison equations have vector-valued solutions of the form
\[
                  h(a)=\sum_{i=1}^r h_i(a_i,a_{i+1}),
                  \qquad a_{r+1}=a_1,
\]
where each $h_i$ is a function of the displayed two binary coordinates.
The adjacent-coordinate representation is what keeps the graph small
despite its exponentially many addresses.

The addressed comparisons encode a monotone circuit computing the
machine's tableau. A private two-dimensional buffer, compared through
its coordinate forms and their sum, makes a singleton source
image force a singleton destination image, while still permitting
full source images. If the machine accepts, a broadcast family forces
all relevant vector sets to be singletons. Buffered identity
self-connections then compare the same coordinate with both buffer
entries, up to their shifts. Adding the entries cancels the coordinate's
two occurrences and expresses the destination
value as a sum of three shifts. This forces the displayed
adjacent-coordinate form.
Every equivalent mate is therefore isomorphic. If the machine does
not accept, two self-connections admit a shift with no global solution
but with the required nonempty image sets, producing a nonisomorphic
equivalent mate.

Section~\ref{sec:wl} establishes the game convention and the upper
bound. Section~\ref{sec:encoding} gives the incidence construction
and recovers all possible equivalent mates.
Section~\ref{sec:offsets} proves the image conditions and global
lifting criterion. Section~\ref{sec:projection} proves the
local-offset extension theorem underlying the converse game strategy.
Sections~\ref{sec:computation} and~\ref{sec:identification} construct
the computation and complete the reduction, with explicit polynomial
size bounds.
